% ============================================================================
% SAMPLE TEXT — /thesis-review v2, section 5.1 `fi-why` (2026-08-21)
% Fills the three placeholder gaps in the opening paragraph (findings V2-2,
% V2-3). SAMPLE ONLY: never paste into chapters/*.tex — adapt in your own
% words. Shown in-conversation the same day. LLM-usage-disclosure artifact.
% ============================================================================

% --- GAP 1: chapter lead-in ("... (Intro, duction into the chapter ...)") ---
% Lands from the eval-framework chapter, announces the build, and sets the
% asymmetry (per the reframe) as what this section contributes.

With the datasets, the evaluation framework, and the explanation machinery in
place, this chapter builds the first model of the lineage: fI-ProtoMF, which
grounds the item side of ProtoMF in catalogue attributes. Before the
construction, this section argues why grounding begins on the item side at
all: the two ingredients of a recommendation turn out to be deeply
asymmetric.

% --- GAP 2: why collaborative embeddings are hard to explain ---
% Replaces "...(makes them hard to explain (why is the collaborative signal
% never explained in itself))". Follows "...often at massive scale,".

The collaborative signal is purely relational: it records who interacted
with what, and nothing about what any of those things are. An embedding
learned from it encodes an item as the crowd of users who chose it, and a
user as the collection of items they chose; no dimension carries a meaning a
human could read off. Whatever explanation such an embedding receives must
be reconstructed after training, which is precisely the post-hoc naming
problem of Section~\ref{sec:XXX-posthoc}.  % <- background-chapter label TBD

% --- GAP 3: the pivot sentence (V2-3) ---
% Replaces "... The two ingredients lens however is more intuitive and will
% therefore inform how the prototypes should be grounded."

The two ingredients of an interaction, in contrast, are far more intuitive
than the collaborative signal itself, and they will therefore inform how the
prototypes should be grounded.
